We begin by the case $\noisestd=0$. As the first $\dimorig$ coordinates are observed exactly, we aim at infering the remaining coordinates of $\state$,%of the random variable,
which correspond to $\bottomstate$.
As such, given an observation $\lmeas$, we aim at sampling from the posterior $\smash{\filter{0}{\lmeas}{\lbottomstate_0}} \propto \bwmarg{0}{\cat{\lmeas}{\lbottomstate_0}}$ with integral form
\begin{equation} 
    \label{eq:posterior:integral_form}
    \filter{0}{\lmeas}{\lbottomstate_0} \propto \textstyle \int \bwmarg{n}{\lstate_n} \left\{ \prod_{s = 1}^{n-1} \bwtrans{}{s}{\lstate_{s+1}}{\lstate_s} \right\} \bwtrans{}{0}{\lstate_1}{\cat{\lmeas}{\lbottomstate_0}} \rmd \lstate_{1:n} \eqsp.
\end{equation}
To solve this problem, we propose to use SMC algorithms \cite{doucet2001sequential, infhidden, chopin2020introduction}, where a set of $N$ random samples, referred to as particles, is iteratively updated to approximate the posterior distribution. The updates involve, at iteration $s$, selecting promising particles from the pool of particles $\pchunk{\particle_{s+1}}{1}{N} = (\particle^1 _{s+1}, \dotsc, \particle^N _{s+1})$ based on a weight function $\uweight{}{s}$, and then apply a Markov transition $\bwopt{s}{\lmeas}{}{}$ to obtain the samples $\pchunk{\particle_{s}}{1}{N}$.
    The transition $\bwopt{s}{\lmeas}{\lstate_{s+1}}{\lstate_s}$ is designed to follow the backward process while guiding the $\dimorig$ top coordinates of the pool of particles $\pchunk{\particle_s}{1}{N}$ towards the measurement $\lmeas$.
%Before we proceed to define the transition kernels, note that under the backward dynamics \eqref{eq:bwker:defn}, $\topstate_t$ and $\bottomstate_t$ are independent conditionally on $\state_{t+1}$ with transition kernels respectively
Note that under the backward dynamics \eqref{eq:bwker:defn}, $\topstate_t$ and $\bottomstate_t$ are independent conditionally on $\state_{t+1}$ with transition kernels respectively
        $\tbwtrans{}{t}{\lstate_{t+1}}{\ltopstate_{t}}  \eqdef \normpdf(\ltopstate_{t}; \tbwmean_{t+1}(\lstate_{t+1}), \sigma^2 _{t+1} \Id_{\dimorig})$ and $
        \bbwtrans{}{t}{\lstate_{t+1}}{\lbottomstate_{t}}  \eqdef \normpdf(\lbottomstate_{t}; \bbwmean_{t+1}(\lstate_{t+1}), \sigma^2 _{t+1} \Id_{\dimx - \dimorig})$
    where $\tbwmean_{t+1}(\lstate_{t+1}) \in \rset^{\dimorig}$ and $\bbwmean_{t+1}(\lstate_{t+1}) \in \rset^{\dimx - \dimorig}$ are such that $\bwmean_{t+1}(\lstate_{t+1}) = \cat{\tbwmean_{t+1}(\lstate_{t+1})}{\bbwmean_{t+1}(\lstate_{t+1})}$ and the above kernels satisfy $\bwtrans{}{t}{\lstate_{s+1}}{\lstate_s} = \tbwtrans{}{t}{\lstate_{t+1}}{\ltopstate_{t}} \bbwtrans{}{t}{\lstate_{t+1}}{\lbottomstate_{t}}$. We consider the following proposal kernels for $t \in [1:n]$,
\begin{equation} 
    \label{eq:optkernel}
    \bwopt{s}{\lmeas}{\lstate_{t+1}}{\lstate_t} \propto \bwtrans{}{t}{\lstate_{t+1}}{\lstate_t} \tfwtrans{t|0}{\lmeas}{\ltopstate_t} \eqsp, \quad \mathrm{where} \quad \tfwtrans{t|0}{\lmeas}{\ltopstate_t} \eqdef \normpdf(\ltopstate_t; \alphacumprod{t}{}^{1/2} \lmeas, (1 - \alphacumprod{t}{}) \Id_{\dimorig}) \eqsp,
\end{equation}
and $\bwopt{n}{\lmeas}{}{}(\lstate_n) \propto \bwmarg{n}{\lstate_n} \tfwtrans{n|0}{\lmeas}{\ltopstate_n}$. For the final step, we define $\bbwopt{0}{\lmeas}{\lstate_1}{\lbottomstate_0} = \bbwtrans{}{0}{\lstate_1}{\lbottomstate_0}$. Using standard Gaussian conjugation formulas, we obtain
\begin{align*}
\bwopt{t}{\lmeas}{\lstate_{t+1}}{\lstate_t} & = \bbwtrans{}{t}{\lstate_{t+1}}{\lbottomstate_t} \cdot \normpdf\left(\ltopstate_t; \mathsf{K}_t \alpha^{1/2} _t \lmeas + (1 - \mathsf{K}_t) \tbwmean_{t+1}(\lstate_{t+1}), (1 - \alphacumprod{t}{}) \mathsf{K}_t \cdot \Id_{\dimorig} \right) \eqsp,\\
\bwmarg{n}{\lstate_n} & = \normpdf(\lbottomstate_n; 0_{\dimx - \dimorig}, \Id_{\dimx - \dimorig})  \cdot \normpdf(\ltopstate_n; \mathsf{K}_n \alphacumprod{n}{}^{1/2} \lmeas, (1 - \alphacumprod{n}{}) \mathsf{K}_n \cdot \Id_{\dimorig})
\end{align*}
where $\mathsf{K}_{t} \eqdef \sigma^2 _{t+1} \big/ (\sigma^2 _{t+1} + 1 - \alpha_t)$. For this procedure to target the posterior $\smash{\filter{0}{\lmeas}{}}$, the weight function $\uweight{}{s}$ is chosen as follows; we set $
    \uweight{}{n-1}(\lstate_n) \eqdef \int \bwtrans{}{n-1}{\lstate_{n}}{\lstate_{n-1}} \tfwtrans{n-1|0}{\lmeas}{\ltopstate_{n-1}} \rmd \lstate_{n-1}
     = \normpdf\left(\alpha_{n-1}^{1/2} \lmeas; \tbwmean_{n}(\lstate_n), \sigma_{n}^2 + 1 - \alpha_s\right)$
and for $\smash{t \in [1:n-2]}$,
\begin{equation}
    \label{eq:weightfunc}
    \uweight{}{t}(\lstate_{t+1}) \eqdef  \frac{\int \tbwtrans{}{t}{\lstate_{t+1}}{\ltopstate_{t}} \tfwtrans{t|0}{\lmeas}{\ltopstate_{t}} \rmd \lstate_{t}}{\tfwtrans{t+1|0}{\lmeas}{\ltopstate_{t+1}}}
    = \frac{\normpdf\left(\alpha_{t}^{1/2} \lmeas; \tbwmean_{t+1}(\lstate_{s+1}), (\sigma_{t+1}^2 + 1 - \alpha_t) \Id_{\dimorig}\right)}{\normpdf\left(\alpha_{t+1}^{1/2}\lmeas; \ltopstate_{t+1}, (1 - \alpha_{t+1}) \Id_{\dimorig} \right)} \eqsp.
\end{equation}
For the final step, we set $\uweight{}{0}(\lstate_1) \eqdef \tbwtrans{}{0}{\ltopstate_1}{\lmeas} \big/ \tfwtrans{1|0}{\lmeas}{\ltopstate_1}$. 
The overall SMC algorithm targeting $\filter{0}{\lmeas}{}$ using the instrumental kernel \eqref{eq:optkernel} and weight function \eqref{eq:weightfunc} is summarized in \Cref{alg:algonoiseless}.
%\vspace{-.3cm}
\begin{algorithm2e}
    \small
    \caption{\algo\ ($\sigma = 0$)}
    \label{alg:algonoiseless} 
    \KwInput{Number of particles $N$}
    \KwOutput{$\pchunk{\particle _0}{1}{N}$}
    \tcp{\scriptsize{Operations involving index $i$ are repeated for each $i \in [1:N]$}}
    $\overline{z}^i _{n} \sim \normpdf(\mathbf{0}_{\dimorig}, \idm_{\dimorig}), \quad \underline{z}^i _{n} \sim \normpdf(\mathbf{0}_{\dimx - \dimorig}, \idm_{\dimx - \dimorig}), \quad \tparticle^i _n = \mathsf{K}_n \alphacumprod{n}{}^{1/2} \lmeas + (1 - \alphacumprod{n}{}) \mathsf{K} _n \overline{z}^i _n , \quad \particle^i _n = \cat{\tparticle^i _n}{\underline{z}^i _n}$;\\
    \For{$s \gets n-1:0$}{
        \If{$s = n-1$}{
            $\uweight{}{n-1}(\particle^i _n) = \normpdf\big(\alphacumprod{n}{}^{1/2} \lmeas; \topvec{\bwmean}_{n}(\particle^i _n), 2 - \alphacumprod{n}{} \big)$;
        }
        \Else{
            $\uweight{}{s}(\particle^i _{s+1}) = \normpdf\left(\alphacumprod{s}{}^{1/2} \lmeas; \topvec{\bwmean}_{s+1}(\particle^i _{s+1}), \sigma_{s+1}^2 + 1 - \alphacumprod{s}{}\right) \big/ \normpdf\left(\alphacumprod{s+1}{}^{1/2}\lmeas; \tparticle^i _{s+1}, 1 - \alphacumprod{s+1}{} \right)$;
        }
        $\anc{i}{s+1} \sim \categorical\big(\{ \uweight{}{s}(\particle^{j} _{s+1}) / \sum_{k = 1}^N \uweight{}{s}(\particle^{k} _{s+1})\} _{j = 1} ^N \big), \quad \overline{z}^i _{s} \sim \normpdf(\mathbf{0}_{\dimorig}, \idm_{\dimorig}), \quad \underline{z}^i _{s} \sim \normpdf(\mathbf{0}_{\dimx - \dimorig}, \idm_{\dimx - \dimorig})$;\\
        $\tparticle^i _s = \mathsf{K}_s \alphacumprod{s}{}^{1/2} \lmeas + (1 - \mathsf{K}_s) \topvec{\bwmean}_{s+1}(\particle^i _{s+1}) + (1 - \alpha_s)^{1/2} \mathsf{K}^{1/2} _s \overline{z}^i _{s}, \quad \bparticle^{i}_{s} = \botvec{\bwmean}_{s+1}(\particle^i _{s+1}) + \sigma _{s+1} \underline{z}^i _{s}$;\\
        Set $\particle^i _s = \cat{\tparticle^i _s}{\bparticle^i _s}$;
    }
    \end{algorithm2e}
We now provide a justification to \Cref{alg:algonoiseless}. Let $\{ \pot{s}{\lmeas}{} \}_{s = 1} ^{n}$ be a sequence of positive functions. Consider the sequence of distributions $\left\{\smash{\filter{s}{\lmeas}{}}\right\}_{s=1}^{n}$ defined as follows; $\filter{n}{\lmeas}{\lstate_n} \propto \bwmarg{n}{\lstate_n} \pot{n}{\lmeas}{\lstate_n}$ and for $t \in [1:n-1]$
\begin{equation}
    \filter{t}{\lmeas}{\lstate_{t}}
     \propto \textstyle \int {\pot{t+1}{\lmeas}{\lstate_{t+1}}}^{-1}{\pot{t}{\lmeas}{\lstate_t}} \bwtrans{}{t}{\lstate_{t+1}}{\lstate_{t}}  \filter{t+1}{\lmeas}{\rmd \lstate_{t+1}} \eqsp,  \label{eq:FKrecursion}
\end{equation}
By construction, the time $t$ marginal \eqref{eq:FKrecursion} is $\filter{t}{\lmeas}{\lstate_t} \propto \bwmarg{t}{\lstate_t} \pot{t}{\lmeas}{\lstate_t}$ for all $t \in [1:n]$.
Then, using $\smash{\filter{1}{\lmeas}{}}$ and \eqref{eq:posterior:integral_form}, we have that $\filter{0}{\lmeas}{\lbottomstate_0} \propto  \int  {\pot{1}{\lmeas}{\lstate_1}}^{-1}{\tbwtrans{}{0}{\ltopstate_1}{\lmeas}} \bbwtrans{}{0}{\lstate_1}{\lbottomstate_0}  \filter{1}{\lmeas}{\rmd \lstate_1}$.

The recursion \eqref{eq:FKrecursion} suggests a way of obtaining a particle approximation of $\filter{0}{\lmeas}{}$; by sequentially approximating each $\filter{t}{\lmeas}{}$ we can effectively derive a particle approximation of the posterior.% using \eqref{eq:time0flow}.
To construct the intermediate particle approximations we use the framework of \emph{auxiliary particle filters} (APF) \cite{pitt1999filtering}.
We focus on the case $\pot{t}{\lmeas}{\lstate_t} = \tfwtrans{t|0}{\lmeas}{\ltopstate_t}$ which corresponds to \Cref{alg:algonoiseless}. The initial particle approximation $\filter{n}{\lmeas}{}$ is obtained by drawing $N$ i.i.d. samples $\pchunk{\particle_{n}}{1}{N}$ from $\bwopt{n}{\lmeas}{}{}$ and setting $\filter{n}{N}{} = N^{-1} \sum_{i = 1}^N \delta_{\particle^{i} _{n}}$ where $\delta_\particle$ is the Dirac mass at $\xi$.
Assume that the empirical approximation of $\filter{t+1}{\lmeas}{}$ is $\filter{t+1}{N}{} = N^{-1} \sum_{i = 1}^N \delta_{\particle^{i} _{t+1}}\eqsp,$ where $\pchunk{\particle_{t+1}}{1}{N}$ are $N$ random variables.
Substituting $\smash{\filter{t+1}{N}{}}$ into the recursion \eqref{eq:FKrecursion} and introducing the instrumental kernel \eqref{eq:optkernel}, we obtain the mixture
\begin{equation} 
\label{eq:particleapprox:intermediate}
    \widehat\phi^N _t (\lstate_t) = \textstyle  \sum_{i = 1}^N {\uweight{}{t}(\particle^{i} _{t+1})}  \bwopt{t}{\lmeas}{\particle^i _{t+1}}{\lstate_t} \big / \sum_{j = 1}^N  \uweight{}{t}(\particle^{j} _{t+1}) \eqsp.
\end{equation}
Then, a particle approximation of \eqref{eq:particleapprox:intermediate} is obtained by  
sampling $N$ conditionally i.i.d. ancestor indices  $\anc{1:N}{t+1} \iid \categorical(\{ \uweight{}{t}(\particle^{i} _{t+1}) / \sum_{j = 1}^N \uweight{}{t}(\particle^{j} _{t+1})\} _{i = 1} ^N)$, 
and then propagating each ancestor particle $\particle^{I^i _{t+1}} _{t+1}$ according to the instrumental kernel \eqref{eq:optkernel}. 
 The final particle approximation is given by
$\filter{0}{N}{} = N^{-1} \sum_{i = 1}^N \delta_{\bparticle^{i} _{0}},$ where $\smash{\bparticle^{i} _{0} \sim \bbwtrans{}{0}{\particle^{\anc{i}{0}} _1}{\cdot}}$, $\smash{\anc{i}{0} \sim \categorical(\{ \uweight{}{0}(\particle^{k} _{1}) \big/ \sum_{j = 1}^N \uweight{}{0}(\particle^j _1) \} _{k = 1} ^N)}$.
 The sequence of distributions $\{ \bwmarg{t}{} \}_{t=0}^{n}$ approximating the marginals of the forward process initialized at $\bwmarg{0}{}$ defines a path that bridges between $\bwmarg{n}{}$ and the prior $\bwmarg{0}{}$ such that the discrepancy between $\bwmarg{t}{}$ and $\bwmarg{t+1}{}$ is small. SMC samplers based on this path are robust to multi-modality and offer an interesting alternative to the geometric and tempering paths traditionally used in the SMC literature, see \cite{dai2022invitation}. Our proposals $\filter{t}{\lmeas}{\lstate_t} \propto \bwmarg{t}{\lstate_t} \tfwtrans{t|0}{\lmeas}{\ltopstate_t}$ inherit the behavior of $\{ \bwmarg{t}{} \}_{t \in \nset}$ and bridge the initial distribution $\filter{n}{\lmeas}{}$ and posterior $\filter{0}{\lmeas}{}$. Indeed, as $\lmeas$ is a noiseless observation of $\state_0 \sim \bwmarg{0}{}$, we may consider $\smash{\alphacumprod{t}{}^{1/2} \lmeas + (1 - \alphacumprod{t}{})^{1/2} \noise_t}$, with $\noise_t \sim \normpdf(\mathbf{0}_{\dimorig}, \Id_{\dimorig})$, as a noisy observation of $X_t \sim \bwmarg{t}{}$ and thus, $\filter{t}{\lmeas}{}$ is the associated posterior.  We illustrate this intuition by considering the following Gaussian mixture (GM) example. We assume that $\bwmarg{0}{\lstate_0} = \sum_{i = 1}^M w_i \cdot \normpdf(\lstate_0; \mu_i, \Id_{\dimx})$ where $M > 1$ and $\{ w_i \}_{i = 1}^M$ are drawn uniformly on the simplex. The marginals of the forward process are available in closed form and are given by $\bwmarg{t}{\lstate_t} = \sum_{i = 1}^M w_i \cdot \normpdf(\lstate_t; \alphacumprod{t}{}^{1/2} \mu_i, \Id_{\dimx})$, which shows that the discrepancy between $\bwmarg{t}{}$ and $\bwmarg{t+1}{}$ is small as long as $\alphacumprod{t}{}^{1/2} - \alphacumprod{t+1}{}^{1/2} \approx 0$. The posteriors $\{ \filter{t}{\lmeas}{} \}_{t \in [0:n]}$ are also available in closed form and displayed in \Cref{fig:bridge_samples}, which illustrates that our choice of potentials ensures that the discrepancy between consecutive posteriors is small.
\begin{figure}
    \centering
    \begin{tabular}{c}
        \begin{tabular}{
            M{0.1125\linewidth}@{\hspace{-0.22\tabcolsep}}
            M{0.1125\linewidth}@{\hspace{-0.22\tabcolsep}}
            M{0.1125\linewidth}@{\hspace{-0.22\tabcolsep}}
            M{0.1125\linewidth}@{\hspace{-0.22\tabcolsep}}
            M{0.1125\linewidth}@{\hspace{-0.22\tabcolsep}}
            M{0.1125\linewidth}@{\hspace{-0.22\tabcolsep}}
            M{0.1125\linewidth}@{\hspace{-0.22\tabcolsep}}
            M{0.1125\linewidth}@{\hspace{-0.22\tabcolsep}}
        }
        \small{$t = 450$} &  \small{$t = 100$} &  \small{$t = 80$} & \small{$t = 70$} & \small{$t=50$} & \small{$t=20$} & \small{$t=15$} & \small{$t=5$} \\
            \includegraphics[width=\sizebridgefig]{images/bridge_distributions/gaussian_mixture_mcgdiff_bridgedistr_step50.pdf}
            & \includegraphics[width=\sizebridgefig]{images/bridge_distributions/gaussian_mixture_mcgdiff_bridgedistr_step400}
            & \includegraphics[width=\sizebridgefig]{images/bridge_distributions/gaussian_mixture_mcgdiff_bridgedistr_step420.pdf}
            & \includegraphics[width=\sizebridgefig]{images/bridge_distributions/gaussian_mixture_mcgdiff_bridgedistr_step430.pdf}
            & \includegraphics[width=\sizebridgefig]{images/bridge_distributions/gaussian_mixture_mcgdiff_bridgedistr_step450.pdf}
            & \includegraphics[width=\sizebridgefig]{images/bridge_distributions/gaussian_mixture_mcgdiff_bridgedistr_step480.pdf}
            & \includegraphics[width=\sizebridgefig]{images/bridge_distributions/gaussian_mixture_mcgdiff_bridgedistr_step485.pdf}
            & \includegraphics[width=\sizebridgefig]{images/bridge_distributions/gaussian_mixture_mcgdiff_bridgedistr_step495.pdf}
        \end{tabular}
    \end{tabular}
    \caption{Display of samples from $\protect\filter{t}{\lmeas}{\lstate_t} \propto \protect\bwmarg{t}{\lstate_t} \protect\tfwtrans{t|0}{\lmeas}{\ltopstate_t}$ for the GM prior. In yellow the samples from $\protect\filter{t}{\lmeas}{}$, in purple those from the prior and in light blue those from the posterior $\protect\filter{0}{\lmeas}{}.$ $n$ is set to $500$.} 
    \label{fig:bridge_samples}
\end{figure}
% The potential $\pot{t}{\lmeas}{\lstate_t} = \tfwtrans{t|0}{\lmeas}{\ltopstate_t}$, and hence Equations \eqref{eq:FKrecursion} and \eqref{eq:time0flow}, is motivated by considering the posterior of the state under the forward process \eqref{eq:forward_def} $\fwdfilter{t}{\lmeas}{\lstate_t} \eqdef  \int \filter{0}{\lmeas}{\lbottomstate_0} \fwtrans{t|0}{\cat{\lmeas}{\lbottomstate_0}}{\lstate_t} \ \rmd \lbottomstate_0$. Indeed, first note that the bridge kernel decomposes across the dimensions as follows,
% \begin{align*} 
%     \revbridge{t-1|t, 0}{\lstate_t,\lstate_0}{\lstate_{t-1}} = \trevbridge{t-1|t, 0}{\ltopstate_t,\ltopstate_0}{\ltopstate_{t-1}} \brevbridge{t-1|t, 0}{\lbottomstate_t,\lbottomstate_0}{\lbottomstate_{t-1}}
% \end{align*}
% where 
% \begin{align*}
%     \trevbridge{t|t+1, 0}{\ltopstate_{t+1},\ltopstate_0}{\ltopstate_{t}} & = \normpdf(\ltopstate_{t}; \bridgemean_t(\ltopstate_0, \ltopstate_{t+1}), \sigma^2 _{t+1} \Id_{\dimorig}) \eqsp,\\ \brevbridge{t-1|t, 0}{\lbottomstate_t,\lbottomstate_0}{\lbottomstate_{t-1}} & = \normpdf(\lbottomstate_{t-1}; \bridgemean_t(\lbottomstate_0, \lbottomstate_t), \sigma^2 _t \Id_{\dimx - \dimorig}) \eqsp,
% \end{align*}
% and $\bridgemean_{t+1}$ is defined in \eqref{eq:bridgeker:mean}. It is then easily seen that 
% \begin{align*} 
%     \tfwtrans{t|0}{\ltopstate_0}{\ltopstate_t} & = \int \trevbridge{t|t+1, 0}{\ltopstate_{t+1},\ltopstate_0}{\ltopstate_{t}} \tfwtrans{t+1|0}{\ltopstate_0}{\rmd \ltopstate_{t+1}} \\
%     \bfwtrans{t|0}{\lbottomstate_0}{\lbottomstate_t} & = \int \brevbridge{t|t+1, 0}{\lbottomstate_{t+1},\lbottomstate_0}{\lbottomstate_{t}} \bfwtrans{t+1|0}{\lbottomstate_0}{\rmd \lbottomstate_{t+1}} \eqsp,
% \end{align*}
% where $\bfwtrans{t|0}{}{}$ is defined analoguously to $\tfwtrans{t|0}{}{}$ in \eqref{eq:optkernel}. Hence, using that $\fwtrans{t|0}{\lstate_0}{\lstate_t} = \tfwtrans{t|0}{\ltopstate_0}{\ltopstate_t} \bfwtrans{t|0}{\lbottomstate_0}{\lbottomstate_t}$, we see that 
% \begin{align*}
%     \fwdfilter{t}{\lmeas}{\lstate_{t}} & = \int \filter{0}{\lmeas}{\rmd \lbottomstate_0} \tfwtrans{t|0}{\lmeas}{\ltopstate_t} \brevbridge{t|t+1, 0}{\lbottomstate_{t+1}, \lbottomstate_0}{\lbottomstate_t} \tfwtrans{t+1|0}{\lbottomstate_0}{\rmd \lbottomstate_{t+1}} \\
%     & = \int \filter{0}{\lmeas}{\rmd \lbottomstate_0} \tfwtrans{t|0}{\lmeas}{\ltopstate_t} \brevbridge{t|t+1, 0}{\lbottomstate_{t+1}, \lbottomstate_0}{\lbottomstate_t} \fwtrans{t+1|0}{\cat{\lmeas}{\lstate_0}}{\rmd \lstate_{t+1}} \eqsp.
% \end{align*}
% Finally, replacing $\lbottomstate_0$ by the vector made of the last $\dimx - \dimorig$ coordinates of its prediction $\xpred{0|s+1}(\lstate_{s+1})$, which we denote by $\bottomxpred{0|s+1}(\lstate_{s+1})$, we see that $\fwdfilter{s}{\lmeas}{}$ satisfies 
% \begin{equation}
%     \label{eq:approximate_recursion}
%     \fwdfilter{s}{\lmeas}{\lstate_s} \approx \int \tfwtrans{s|0}{\lmeas}{\ltopstate_s} \bbwtrans{}{s}{\lstate_{s+1}}{\lbottomstate_s} \fwdfilter{s+1}{\lmeas}{\rmd \lstate_{s+1}} \eqsp.
% \end{equation}
% According to this recursion, the transition from $\lbottomstate_{s}$ should follow the backward process conditioned by $\lstate_{s+1}$, while that from $\ltopstate_{s}$ should follow the forward one conditioned on the measurement $\lmeas$ at time $0$. \Cref{eq:FKrecursion} reflects this behaviour, while \eqref{eq:time0flow} ensures that the procedure ultimately produces the posterior as the final marginal.
The idea of using the forward diffused observation to guide the observed part of the state, as we do here through $\tfwtrans{t}{\lmeas}{\ltopstate_t}$, has been exploited in prior works but in a different way. For instance, in \cite{song2020score,song2022solving} the observed part of the state is directly replaced by the forward noisy observation and, as it has been noted \cite{trippe2023diffusion}, this introduces an irreducible bias. Instead, \algo\ weights the backward process by the density of the forward one conditioned on $\lmeas$, resulting in a natural and consistent algorithm.

We now establish the convergence of \algo\ with a general sequence of \emph{potentials} $\{ \pot{s}{\lmeas}{} \}_{s = 1} ^n$. We consider the following assumption on the sequence of potentials $\{ \pot{t}{\lmeas}{} \}_{t = 1} ^n$. 
\begin{hypA} 
    \label{assp:bounded_weights}
$\displaystyle \sup_{\lstate \in \R^{\dimx}} \tbwtrans{}{0}{\lstate}{\lmeas} / \pot{1}{\lmeas}{\lstate} < \infty$ and $\displaystyle\sup_{\lstate \in \R^{\dimx}} \int \pot{t}{\lmeas}{\lstate_{t}} \bwtrans{}{t}{\lstate}{\lstate_{t}} \rmd \lstate_{t} \big/ \pot{t+1}{\lmeas}{\lstate} < \infty$ for all $t \in [1:n-1]$.
\end{hypA}
The following exponential deviation inequality is standard and is a direct application of \cite[Theorem 10.17]{douc2014nonlinear}. In particular, it implies a $\bigo(1 / \sqrt{N})$ bound on the mean squared error $\| \filter{0}{N}{h} - \filter{0}{\lmeas}{h} \|_2$.
\begin{proposition} Assume \assp{\ref{assp:bounded_weights}}. There exist constants $c_{1,n}, c_{2,n} \in (0, \infty)$ such that, for all $N \in \nset$, $\varepsilon > 0$ and bounded function $h : \rset^{\dimx} \mapsto \rset$,
$
    \pP \left[ \big| \filter{0}{N}{h} - \filter{0}{\ltrmeas}{h} \big| \geq \varepsilon \right] \leq c_{1,n} \exp(- c_{2,n} N \varepsilon^2 \big/ |h|^2 _\infty )
$
where $|h|_\infty \eqdef \sup_{x \in \rset^{\dimx}} |h(x)|$.
\end{proposition}
We also furnish our estimator with an explicit non-asymptotic bound on its bias. Define $\smash{\bfilter{0}{N}{} = \pE \big[ \filter{0}{N}{} ]}$ where $\filter{0}{N}{} = N^{-1} \sum_{i = 1}^N \delta_{\bparticle^i _0}$ is the particle approximation produced by \Cref{alg:algonoiseless} and the expectation is with respect to the law of $(\particle^{1:N}_{1:n}, \anc{1:N}{0:n-1})$.
Define for all $t \in [1:n]$, $\filter{t}{\star}{\lstate_t} \propto \bwmarg{t}{\lstate_t} \int \delta_\lmeas(\rmd \ltopstate_0) \bwtrans{}{0|t}{\lstate_t}{\lstate_0} \rmd \lbottomstate_0 \eqsp,$ where $\bwtrans{}{0|t}{\lstate_t}{\lstate_0} \eqdef \int \left\{ \prod_{s = 0} ^{t-1} \bwtrans{}{s}{\lstate_{s+1}}{\lstate_s} \right\} \rmd \lstate_{1:t-1}$.
\begin{proposition} \label{lem:kldiv:filtering}
   It holds that
    \begin{equation}
        \kldivergence{\filter{0}\lmeas{}}{\bfilter{0}{N}{}} \leq \mathsf{C}^{\lmeas} _{0:n}(N-1)^{-1} + \mathsf{D}^{\lmeas} _{0:n} N^{-2} \eqsp,
    \end{equation}
    where  $\mathsf{D}^\lmeas _{0:n} > 0$,
    $\mathsf{C}^\lmeas _{0:n} \eqdef 
            \textstyle\sum_{t = 1} ^n \int \frac{\normconst_t / \normconst_0}{\pot{t}{\lmeas}{z_t}} \left\{ \int \delta_{\lmeas}(\rmd \ltopstate_0)  \bwtrans{}{0|t}{z_t}{x_0} \rmd \lbottomstate_0 \right\} \filter{t}{\star}{\rmd z_t} $
    and $\normconst_t \eqdef \int \pot{t}{\lmeas}{\lstate_t} \bwmarg{t}{\rmd \lstate_t}$ for all $t \in [1:n]$ and $\normconst_0 \eqdef \int \delta_\lmeas(\rmd \ltopstate_0) \bwmarg{0}{\lstate_0} \rmd \lbottomstate_0$. If furthermore \assp{\ref{assp:bounded_weights}} holds then both $\mathsf{C}^{\lmeas} _{0:n}$ and $\mathsf{D}^{\lmeas} _{0:n}$ are finite.
\end{proposition}
The proof of \Cref{lem:kldiv:filtering} is postponed to \Cref{proof:kldiv:filtering}. \assp{\ref{assp:bounded_weights}} is an assumption on the equivalent of the weights $\{\uweight{}{t}\}_{t = 0}^n$ with a general sequence of potentials $\{ \pot{t}{\lmeas}{}\}_{t = 1} ^n$ and is not restrictive as it can be satisfied by setting for example $\pot{s}{\lmeas}{\lstate_s} = \tfwtrans{s|0}{\lmeas}{\ltopstate_s} + \delta$ where $\delta > 0$. The resulting algorithm is then only a slight modification of the one described above, see \Cref{sec:proof_lems} for more details. 
It is also worth noting that \Cref{lem:kldiv:filtering} combined with Pinsker's inequality implies that the bias of \algo\ goes to $0$ with the number of particle samples $N$ for fixed $n$. We have chosen to present a bound in Kullback--Leibler (KL) divergence, inspired by \cite{andrieu2018uniform, huggins2019sequential}, as it allows an explicit dependence on the modeling choice $\{ \pot{s}{\lmeas}{} \}_{s = 1} ^n$, see \Cref{lem:closed_form_bound}.
Finally, unlike the theoretical guarantees established for \smcdiff\ in \cite{trippe2023diffusion}, proving the asymptotic exactness of our methodology w.r.t. to the generative model posterior does not require having $\bwmarg{s+1}{\lstate_{s+1}} \bwtrans{}{s}{\lstate_{s+1}}{\lstate_{s}} = \bwmarg{s}{\lstate_s} \fwtrans{s+1}{\lstate_s}{\lstate_{s+1}}$ for all $\smash{s \in [0:n-1]}$, %i.e. that the one step forward kernel is the time reversal of the backward one,
which does not in hold in practice. %As such, \smcdiff\ exhibits a non-vanishing asymptotic bias.